\scanpage{138}
% archive foliation: 137
Soit $f:X\mathbin{\text{---}}S$ un morphisme se factorisant en
\diagram{\begin{tikzcd}X\arrow[r,dash,"\varphi"]&\Sigma\arrow[r,dash,"\psi"]&S\end{tikzcd}}
, avec $\psi$ fini, et $\varphi$ tel que $\uline{O}_{\Sigma}\simeq\varphi_*(O_X)$.

1) Alors on a
\diagram{\begin{tikzcd}\bullet\arrow[r]&\mathrm{Pic}(\Sigma)\arrow[r,dash]&\mathrm{Pic}(X)\arrow[r,dash]&\mathrm{Pic}'(X/S)\end{tikzcd}}
[suite exacte provenant de la suite spectrale de Leray].

\deleted{\unsure{et}} \uline{Si} $X/\Sigma$ admet une section, le dernier hom est un isom, i.e. dans ce cas
\[
\mathrm{Pic}'(X/S)\simeq\mathrm{Pic}(\unsure{X})/\mathrm{Pic}(\Sigma)
\]

Ce dernier fait se voit ainsi: \deleted{\unsure{Pour}} Soit $g:\Sigma\longrightarrow X$ une section de $X/\Sigma$; alors

i) \deleted{\unsure{pour}}\inserted{un} Module inversible $g$-rigidifié sur $X$ n'a pas d'aut. non trivial, ii) Tout Module inversible $\uline{L}$ sur $X$ est \deleted{\unsure{de norme}} $\mathrm{Pic}'(X/S)$ \uline{équivalent} à \deleted{l'image} un Module inversible $g$-rigidifié \inserted{[savoir $\uline{L}\otimes\varphi^*(g^*(\uline{L}))$]} [déterminé à isom unique près].

[On utilise le fait que $\mathrm{Pic}'(\Sigma/S)=0$].

On suppose pour la suite $f$ qu. cpt et qu. sép

